\chapter{Settlement Risk}\label{ch:m2:settlement-risk}

On 26 June 1974 the German banking supervisor closed Bankhaus Herstatt, a
German bank that had lost heavily on speculative currency positions. It did so
in the middle of the German business day, before the markets in New York had
opened. By then Herstatt had received, through the German payment system, the
Deutsche marks it had bought two days earlier; it had not yet paid the dollars
it had sold, which were due later that day in New York. Its counterparties had
paid and would not be paid. Several banks were hit, and the main American
clearing system for dollar payments closed for a day. Until then, a currency
trader's risk had been thought to be the price; Herstatt showed that the whole
principal of a trade could be lost between two payments in two time zones
(\cref{fig:m2:settlement-risk:herstatt}). Half
a century later, the central banks' survey of April 2025 found that about USD~1.4
trillion a day was still settled in a way fully exposed to that risk. This
chapter explains \omterm{def:m2:settlement-risk:risk}{settlement risk}, the payment-versus-payment system built to
remove it, the netting that reduces it, and what remains.

\section{Herstatt}

\begin{definition}[Settlement risk, Herstatt risk]\label{def:m2:settlement-risk:risk}
The \emph{settlement risk}\index{settlement risk} of a foreign exchange trade is
the risk that a party pays the currency it sold and does not receive the currency
it bought. In its time-zone form, the payment of one leg becoming irrevocable
hours before the other leg is due, it is called \emph{Herstatt
risk}\index{Herstatt risk}.
\end{definition}

The risk is not the market move, which is small over a day, but the principal.
A bank that sells yen for dollars pays its yen in the Asian morning and receives
its dollars in the American afternoon; if its counterparty fails in between, it
has lost the yen and does not get the dollars. The exposure lasts from the
moment its own payment can no longer be recalled until the moment it knows the
other payment has arrived.

\begin{definition}[Unilateral irrevocable payment time]\label{def:m2:settlement-risk:irrevocable}
The \emph{unilateral irrevocable payment time}\index{unilateral irrevocable payment time}
of a currency, for a given bank, is the deadline after which it can no longer
cancel a payment it has instructed in that currency. A trade's settlement
exposure runs from the irrevocable time of the currency sold to the confirmed,
final receipt of the currency bought, and may be extended by the time it takes
to check the receipt and by a failed payment.
\end{definition}

\begin{omfigure}
\centering
\begin{tikzpicture}[font=\small,x=0.42cm]
\draw[thick,-{Stealth}] (-3,0) -- (25,0) node[right,font=\footnotesize] {hours, GMT};
\foreach \h/\l in {-2/{22:00},4/{4:00},10/{10:00},16/{16:00},22/{22:00}} {\draw (\h,0.1) -- (\h,-0.1) node[below,font=\scriptsize] {\l};}
\fill[omDef!25] (-2,0.5) rectangle (7,1.0); \node[font=\footnotesize,anchor=west] at (-2,1.25) {yen payments (Tokyo)};
\fill[omThm!25] (6,1.8) rectangle (16,2.3); \node[font=\footnotesize,anchor=west] at (6,2.55) {euro payments};
\fill[omProp!25] (13,3.1) rectangle (21,3.6); \node[font=\footnotesize,anchor=west] at (13,3.85) {dollar payments (New York)};
\draw[omProp,very thick,{Stealth}-{Stealth}] (-2,4.7) -- (21,4.7);
\node[omProp,font=\footnotesize,anchor=south] at (9.5,4.75) {a bank that sells yen for dollars is exposed for about 23 hours};
\end{tikzpicture}

\omcaption{Why the time zones create the risk. Payments in each currency are made
during its own system's hours; a bank that sells yen and buys dollars cannot
recall its yen after the Tokyo deadline and does not know it has its dollars until
New York's afternoon. Illustrative hours, GMT, on the settlement
day.}\label{fig:m2:settlement-risk:zones}
\end{omfigure}

\begin{omfigure}
\centering
\begin{tikzpicture}[font=\small,x=0.5cm]
\draw[thick,-{Stealth}] (0,0) -- (26,0) node[right,font=\footnotesize] {time};
\draw (2,0.1) -- (2,-0.1) node[below,font=\footnotesize,align=center] {24 June\\trades};
\draw (12,0.1) -- (12,-0.1) node[below,font=\footnotesize,align=center] {26 June, midday\\Germany: bank closed};
\draw (20,0.1) -- (20,-0.1) node[below,font=\footnotesize,align=center] {New York opens:\\dollars due};
\draw[omDef,thick,-{Stealth}] (8,1.9) -- (11.8,0.25);
\node[omDef,font=\footnotesize,align=center] at (8,2.4) {marks received by Herstatt\\through the German system};
\draw[omProp,thick,-{Stealth}] (22,1.9) -- (20.2,0.25);
\node[omProp,font=\footnotesize,align=center] at (22.5,2.4) {dollars never paid\\to its counterparties};
\end{tikzpicture}

\omcaption{The Herstatt day, 26 June 1974. Counterparties had paid their marks for
trades done two days earlier; the bank was closed in the middle of the German
business day, before New York opened, and the dollars it owed were not paid.
After the ECB's account of 2007. Schematic.}\label{fig:m2:settlement-risk:herstatt}
\end{omfigure}

The central banks of the G10 answered in 1996 with a strategy on three tracks:
banks should measure and control their settlement exposures like any credit
exposure; industry groups should build services that settle both legs together;
central banks should push both and adapt their payment systems. The industry's
answer was CLS, which began operating in September 2002.

\section{Payment versus payment}

\begin{definition}[Payment versus payment]\label{def:m2:settlement-risk:pvp}
\emph{Payment versus payment}\index{payment versus payment} (PvP) is a settlement
mechanism in which the final payment of one currency occurs if, and only if, the
final payment of the other currency occurs.
\end{definition}

PvP is to currencies what delivery versus payment is to securities (One Quant
Book 1, chapter 5). CLS operates it for 18 of the main currencies: it settles
each trade's two legs together, or not at all. \omterm{def:m2:settlement-risk:risk}{Settlement risk} disappears; in its
place each member must pay in, on time, the currencies it owes, a liquidity demand
that is smaller when many trades are settled against one another.

\begin{proposition}[What PvP removes and what it leaves]\label{prop:m2:settlement-risk:pvp}
Under PvP a failed counterparty's trades are not settled and the survivor keeps
the currency it would have paid: its loss is the cost of replacing the trade at
the new market rate, the replacement risk, not the principal. Replacement risk is
of the order of the day's exchange-rate move, a small fraction of principal.
\end{proposition}

\begin{proof}
If the counterparty's leg is not paid, the survivor's leg is not paid either, so
no principal leaves it. It must still buy the currency it needed elsewhere, at the
rate of that moment instead of the agreed one: the difference is its loss.
\end{proof}

\section{Netting}

\begin{definition}[Payment netting]\label{def:m2:settlement-risk:netting}
\emph{Payment netting}\index{payment netting} is the settlement, between two
parties and for each currency and value date, of only the net of the amounts they
owe each other, in place of every trade's gross payment.
\end{definition}

A bank that sells EUR~400 million for dollars to a counterparty and buys back
EUR~300 million from it on the same value date need pay only EUR~100 million net;
its exposure falls from the full 400 million to 100 million. The BIS estimated
that pre-settlement netting removed about USD~1.3 trillion a day of payments in
April 2022. Netting reduces the amount at risk but not its nature: what is still
paid gross can still be lost.

\begin{example}[One bank's settlement day]\label{ex:m2:settlement-risk:book}
A bank's book for one value date has eight trades with four counterparties,
USD~1.53 billion in all: with A, it sells EUR~400 million and buys EUR~300 million
against dollars; with B, it sells yen for USD~250 million and buys yen for
USD~150 million; with C, it sells euros for sterling (200 million) and dollars for
euros (100 million); with D, it sells an emerging-market currency outside CLS for
USD~80 million and buys it back for USD~50 million. Settled gross with the
illustrative hours of \cref{fig:m2:settlement-risk:zones}, its exposure peaks at
USD~1.33 billion in the afternoon, when the dollar payments are due and the euro
ones have gone. Netted per counterparty and currency, the payments fall to USD~430
million and the peak to USD~430 million; settled through PvP in every currency but
the last, the peak is USD~30 million, the net of D's trades
(\cref{fig:m2:settlement-risk:exposure}).
\end{example}

\begin{omfigure}
\begin{tikzpicture}
\begin{axis}[width=11cm,height=5cm,
  xlabel={hour of the settlement day (GMT)},ylabel={exposure (USD million)},
  tick label style={font=\small},label style={font=\small},
  xmin=-4,xmax=24,ymin=0,ymax=1450,grid=major,grid style={gray!25},
  legend columns=3,legend style={font=\footnotesize,at={(0.5,-0.3)},anchor=north,draw=none,fill=none,/tikz/every even column/.append style={column sep=8pt}},
  table/col sep=comma]
\addplot[omProp,very thick,const plot] table[x=t,y=gross]{figdata/markets-2/20-settlement-risk/exposure.csv};
\addplot[omThm,very thick,dashed,const plot] table[x=t,y=net]{figdata/markets-2/20-settlement-risk/exposure.csv};
\addplot[omDef,very thick,dotted,const plot] table[x=t,y=pvp]{figdata/markets-2/20-settlement-risk/exposure.csv};
\legend{gross,netted,netted and PvP}
\end{axis}
\end{tikzpicture}

\omcaption{The settlement exposure of the book of
\cref{ex:m2:settlement-risk:book} through the settlement day. Gross settlement
exposes the bank to most of its day's purchases at once in the afternoon; netting
cuts the exposure by two thirds; PvP leaves only the currency it does not cover.
Illustrative; data: the chapter's tutorial.}\label{fig:m2:settlement-risk:exposure}
\end{omfigure}

\section{What remains unprotected}

\omterm{def:m2:settlement-risk:risk}{Settlement risk} has not gone. In April 2022, by the BIS's estimate, USD~2.2 trillion
of the day's deliverable FX turnover, almost a third, was still settled at risk; in
the largest trading centres 20--40\% of turnover, in some smaller ones more than
three quarters. The 2025 survey measured it differently, by settlement method: just
over a third of the day's settlement, USD~5.2 trillion, went through PvP; USD~7.6
trillion used methods that mitigate the risk without removing it, such as netting,
settlement within a group, or bank accounts with timing controls; and USD~1.4
trillion, 10\%, was settled gross between the two parties. The main reasons given
were that the counterparty had no access to PvP or that the currency was not
supported (\cref{fig:m2:settlement-risk:methods}).

\begin{omfigure}
\begin{tikzpicture}
\begin{axis}[width=11cm,height=4.6cm,xbar,bar width=12pt,
  xlabel={USD trillion a day, April 2025},
  tick label style={font=\small},label style={font=\small},
  xmin=0,xmax=9.2,ymin=-0.6,ymax=2.6,grid=major,grid style={gray!25},
  ytick={0,1,2},yticklabels={gross bilateral,{netting, intragroup, timing controls},payment versus payment},
  yticklabel style={font=\footnotesize,text width=4.2cm,align=right},
  nodes near coords,every node near coord/.append style={font=\scriptsize},
  table/col sep=comma]
\addplot[fill=omThm!60,draw=omThm] coordinates {(1.4,0) (7.6,1) (5.2,2)};
\end{axis}
\end{tikzpicture}

\omcaption{How the world's FX settlement was done in April 2025, by method: just over
a third by \omterm{def:m2:settlement-risk:pvp}{payment versus payment}, more than half by methods that reduce but do not
remove \omterm{def:m2:settlement-risk:risk}{settlement risk}, and a tenth, USD~1.4 trillion a day, gross between the two
parties. Data: BIS Triennial Survey, as reported in the \emph{BIS Quarterly
Review}, June 2026.}\label{fig:m2:settlement-risk:methods}
\end{omfigure}

Measuring one's own exposure is itself a discipline. A bank's exposure to a trade
does not end when the currency it bought is due, but when it knows it has arrived;
the ECB described, in 2007, the irrevocable period that ends with final receipt, an
uncertain period while the bank has not yet checked it, and a failed period if the
payment did not come. A bank that reconciles its incoming payments only the next
morning carries its whole purchases overnight in the uncertain period, and a failed
payment extends the exposure until it is resolved.

The losses are not only history. The BIS article of 2022 recalled KfW's loss of
EUR~300 million when Lehman Brothers failed in 2008 and Barclays's loss of USD~130
million to a small currency exchange in March 2020. The pattern is Herstatt's: a
payment made on the assumption that the other would follow.

\begin{dated}{2026-09}{Settlement today}\label{dat:m2:settlement-risk:today}
BIS Triennial Survey, April 2025: USD~5.2 trillion a day settled by PvP, USD~7.6
trillion by methods that mitigate the risk, USD~1.4 trillion gross bilateral (10\%).
CLS settles 18 currencies, on average more than USD~8 trillion a day, for more than
75 settlement members.
\end{dated}

\section{Tutorial: the settlement-exposure timeline}\label{tut:m2:settlement-risk:timeline}

\begin{tutorial}
\textbf{Goal.} Compute a book's settlement exposure through the day, gross, netted
and with PvP, from each currency's cancellation deadline and confirmation time.
\textbf{End state:} \cref{fig:m2:settlement-risk:exposure},
\cref{ex:m2:settlement-risk:book} and the numbers of the weekend problem.
\begin{tutsteps}
\item \textbf{The exposure window} of a trade and the exposure at a time.
\omcode{code/firm/settlerisk/firm_settlerisk.py}{20}{28}{Exposure window of a trade and a book's exposure at a time.}\label{lst:m2:settlement-risk:window}
\item \textbf{Netting and the day's profile}, with PvP currencies removed.
\omcode{code/firm/settlerisk/firm_settlerisk.py}{31}{61}{Payment netting per counterparty, and the exposure profile.}\label{lst:m2:settlement-risk:profile}
\item \textbf{Run} \texttt{settle\_demo.summary()} and \texttt{fig\_settle.py}.
\end{tutsteps}
\textbf{What to change next.} Add an uncertain period, the hours before the bank
confirms a receipt, to each currency, and a failed payment that extends a trade's
exposure to the next day; see how the peak moves.
\end{tutorial}

\section{Build: the settlement-exposure calculator}\label{bld:m2:settlement-risk:settlerisk}

\begin{build}
\bfield{Purpose} The miniature firm settles its own FX trades and those of its
clients: it must know, hour by hour, how much it could lose if a counterparty
failed, set limits on it, and see what netting and PvP would save.
\bfield{Interface} \texttt{Trade(counterparty, sell, buy, value\_usd)};
\texttt{window(trade, times)}; \texttt{exposure\_at(t, trades, times)};
\texttt{net\_by\_counterparty(trades)}; \texttt{profile(trades, times, pvp)};
\texttt{peak(profile)}.
\bfield{Rules} Times per currency as (cancellation deadline, confirmation) in GMT
hours of the settlement day; values in dollars; netting per counterparty and
currency; trades between two PvP currencies carry no settlement exposure.
\bfield{Acceptance tests} \texttt{code/firm/settlerisk/tests/}: the yen-for-dollar
window and the empty reverse window; exposure summed and zero after the last
confirmation; PvP removes it; netting reduces payments and the peak.
\bfield{Stretch} Per-counterparty limits and alerts; uncertain and failed periods;
PvP liquidity: the pay-ins a member needs by currency and hour.
\end{build}

\begin{omsources}
\item European Central Bank, \emph{Financial Stability Review}, December 2007,
Box 19, ``More than thirty years after the Herstatt case''.
\item M.~Glowka and T.~Nilsson, ``FX settlement risk: an unsettled issue'',
\emph{BIS Quarterly Review}, December 2022.
\item M.~Drehmann, P.~McGuire, T.~Shirakami, M.~Conway and N.~Lovell,
``Uncovering FX settlement risk: new measures from the 2025 BIS Triennial Survey'',
\emph{BIS Quarterly Review}, June 2026.
\item CLS Group, company information.
\end{omsources}

\section{Exercises}

\begin{exercise}[$\star$]\label{exo:m2:settlement-risk:1}
With the illustrative hours of \cref{fig:m2:settlement-risk:zones}, how long is a
bank exposed when it sells yen for dollars, and when it sells dollars for yen?
\end{exercise}

\begin{exercise}[$\star$]\label{exo:m2:settlement-risk:2}
What does a bank lose if its counterparty fails during the exposure window, under
gross settlement and under PvP?
\end{exercise}

\begin{exercise}[$\star$]\label{exo:m2:settlement-risk:3}
Two banks have ten trades in EURUSD for the same value date: the first sells EUR~1
billion in total and buys EUR~900 million. What is paid under \omterm{def:m2:settlement-risk:netting}{payment netting}?
\end{exercise}

\begin{exercise}[$\star\star$]\label{exo:m2:settlement-risk:4}
In \cref{ex:m2:settlement-risk:book}, give the net payments with each counterparty.
\end{exercise}

\begin{exercise}[$\star\star$]\label{exo:m2:settlement-risk:5}
Why does PvP create a liquidity need, and how does settling many trades together
reduce it?
\end{exercise}

\begin{exercise}[$\star\star$]\label{exo:m2:settlement-risk:6}
Why is \omterm{def:m2:settlement-risk:risk}{settlement risk} higher in emerging-market currencies and with smaller
counterparties?
\end{exercise}

\begin{exercise}[$\star\star\star$]\label{exo:m2:settlement-risk:7}
\emph{Coding.} With \texttt{profile}, give the exposure of the book at 05:00,
08:00, 14:00 and 17:00, gross and netted.
\end{exercise}

\begin{exercise}[$\star\star\star$]\label{exo:m2:settlement-risk:8}
\emph{Find the flaw.} ``Our FX exposure to a counterparty is the market value of our
open forwards with it, a few million; \omterm{def:m2:settlement-risk:risk}{settlement risk} is negligible.'' Correct it.
\end{exercise}

\section{Problem: The Unsettled Trillion}

\begin{problem}[{Weekend problem --- a bank's settlement day}]\label{pb:m2:settlement-risk:1}
A bank has the book of \cref{ex:m2:settlement-risk:book} for tomorrow's value date,
with the illustrative payment hours of \cref{fig:m2:settlement-risk:zones} (yen
from 22:00 the evening before to 07:00, euro and sterling 06:00--16:00 and
07:00--16:00, dollar 13:00--21:00, the emerging currency 04:00--12:00).

\medskip\noindent\textbf{Part I --- Gross.}
\begin{enumerate}
\item Give the day's gross purchases and the peak exposure.
\item When is the peak, and why then?
\item Which single trade is exposed longest, and for how long?
\item What is the exposure at 05:00?
\item If counterparty A failed at 14:00, what would the bank lose?
\end{enumerate}

\medskip\noindent\textbf{Part II --- Netting.}
\begin{enumerate}[resume]
\item Give the net payments with each counterparty.
\item Give the netted peak.
\item What would the bank lose if A failed at 14:00 after netting?
\item What does netting need legally to hold in a failure?
\item What remains with D after netting, and why?
\end{enumerate}

\medskip\noindent\textbf{Part III --- PvP.}
\begin{enumerate}[resume]
\item Give the peak with PvP for every currency but the emerging one.
\item What does the bank now risk with A, B and C?
\item What does PvP cost the bank?
\item How could the last USD~30 million be protected?
\item Why do many banks still settle some trades gross?
\end{enumerate}

\medskip\noindent\textbf{Part IV --- Judgement.}
\begin{enumerate}[resume]
\item How should a bank set a settlement limit for a counterparty?
\item What did the Herstatt case change in how banks see FX risk?
\item What does the 2025 survey's 10\% mean for the market as a whole?
\item State the \emph{named result}: the book's peak exposure gross, netted and
with PvP.
\item In one sentence: what is \omterm{def:m2:settlement-risk:risk}{settlement risk}, and what removes it?
\end{enumerate}
\end{problem}

\section{Interview questions}

\begin{interviewq}[$\star$ \iqroles{bank, developer}]\label{iq:m2:settlement-risk:1}
What is \omterm{def:m2:settlement-risk:risk}{Herstatt risk}, and how does CLS remove it?
\end{interviewq}

\begin{interviewq}[$\star$ \iqroles{bank}]\label{iq:m2:settlement-risk:2}
What is the difference between \omterm{def:m2:settlement-risk:risk}{settlement risk} and replacement risk?
\end{interviewq}

\begin{interviewq}[$\star\star$ \iqroles{developer}]\label{iq:m2:settlement-risk:3}
How would you compute a bank's intraday settlement exposure to each counterparty?
\end{interviewq}

\begin{interviewq}[$\star\star$ \iqroles{bank, trader}]\label{iq:m2:settlement-risk:4}
Why is not all FX settled through PvP?
\end{interviewq}

\begin{interviewq}[$\star\star$ \iqroles{researcher, bank}]\label{iq:m2:settlement-risk:5}
How much can \omterm{def:m2:settlement-risk:netting}{payment netting} reduce settlement exposure, and what limits it?
\end{interviewq}

\begin{interviewq}[$\star\star\star$ \iqroles{developer, bank}]\label{iq:m2:settlement-risk:6}
Design a system that blocks new trades with a counterparty when its projected
settlement exposure on a value date would exceed its limit.
\end{interviewq}
